% !TEX root = ../main.tex
%-----------------------------------------------------------------------
%  The five findings, the two figures, and the companion animations.
%  \input by ../main.tex -- not compiled on its own.
%-----------------------------------------------------------------------
\hd{Finding 1 --- the transport responds strongly, and the signal is real.}
The cost $J$ rises steadily with mixing, from 0.214 at $0.25\times$ to 0.510
at $32\times$. To judge whether such differences could be chance, the same
transport index was reconstructed over the entire 200-year spin-up from its
2{,}402 saved restart files, giving the natural year-to-year variability of
$J$: about $\pm0.007$ (annual means) to $\pm0.012$ (individual samples).
Every member's departure from the reference exceeds this by at least
$2.4\sigma$, and most by far more. One subtlety: the reference sits about
$0.07$ \emph{above} the trend through the members. This is not an error ---
perturbing the mixing in \emph{either} direction weakens the transport
temporarily while the ocean adjusts, so all seven adjusted members sit below
the never-perturbed reference. A $1\times$ control member (rerunning the
ten-year leg with unchanged mixing) is the clean way to confirm this and is
the first recommended follow-up.

\hd{Finding 2 --- where the adjoint stays healthy, the maps keep their shape
but not their size.}
For the stable members, the five-year accumulated temperature-sensitivity
map correlates with the reference at 0.82 ($0.5\times$) and 0.85 ($2\times$),
with roughly 80\,\% of grid points keeping their sign --- the geography of
``where the transport is sensitive'' survives. The amplitude does not: overall
magnitudes shift to $1.37\times$ and $0.77\times$ the reference. The
sensitivity to \kv{} itself is markedly more fragile: its pattern correlation
is only 0.21 already at $2\times$, and the globally summed
$\partial J/\partial\kv$ runs $-251 \to -436 \to -78 \to +673$ across
$0.5\times$--$16\times$ --- it changes sign within the ensemble. The one
field a mixing-aware surrogate most needs to predict is the one that varies
fastest with mixing.

\hd{Finding 3 --- four of seven adjoint computations go unstable.}
The forward simulations and cost values are healthy in every member; what
fails is the backward computation, which amplifies small signals
explosively once it runs far enough into the past around a sufficiently
altered ocean state. The instability grows in short bursts
($e$-folding 1--3 days), always ignites in the Southern Ocean channel
(38--51\dg S, mostly at depth, far from the 26\dg N section), and its onset is
\emph{not} ordered by perturbation size: $32\times$ fails first
(0.36\,yr back), then $4\times$, $8\times$, and even the mild $0.25\times$
(2.6\,yr back), while $0.5\times$, $2\times$ and $16\times$ never fail.
Stability is a property of each member's own adjusted ocean state, so it
cannot be guaranteed by simply capping the perturbation. Figure~\ref{fig:blowup}
shows the $8\times$ case developing. All runs used the plain adjoint; the
configuration already carries a stabilised variant (\code{adjViscBoost},
which smooths the backward computation) that is the natural remedy to test.

\begin{figure}[t]
\centering
\includegraphics[width=\textwidth]{fig_m5_blowup_montage}
\caption{\small The instability in space: the $8\times$ member's
temperature-sensitivity magnitude (logarithmic colour scale) at 231\,m depth,
read left to right and top to bottom as the adjoint works backward from the
cost period. The map is well behaved one year back (top left panels), a
localised anomaly ignites in the Southern Ocean near 1.8 years back, and
within about 50 model days the instability fills the domain (bright panels).
The green line marks the 26\dg N section that defines the cost. An
interactive version with a time slider accompanies this brief
(\code{ADJtheta\_M5\_blowup\_log10.html}).}
\label{fig:blowup}
\end{figure}

\hd{Finding 4 --- the straight-line prediction fails at every tested change.}
The adjoint's gradients predict, to first order, how much $J$ should change
for a given change in \kv{} and in the adjusted initial state. Comparing that
prediction against the measured changes: four of seven members get the wrong
sign, and none is predicted within 30\,\%. The linear approximation is
therefore trustworthy only for \kv{} changes smaller than a factor of two ---
the smallest step tested. The failure comes mainly from the ten-year ocean
adjustment (its temperature and salinity contributions are individually large
and nearly cancel), not from the direct effect of \kv{} within the five-year
window.

\hd{Finding 5 --- two incidental discoveries about the cost function.}
Working through the comparison exposed two properties of the compiled cost
worth recording. (i)~$J$ averages the transport only over the \emph{final 30
days} of the five-year window (a package default we had not overridden), so
direct cost information enters the adjoint only in the last month ---
everything earlier is dynamics. (ii)~$J$'s normalisation is computed
per processor tile, so its numerical value depends on how the domain is split
across processors; all runs here share one decomposition, so comparisons
stand, but a rebuild on different processor counts would change $J$ itself.
Both are now documented in the working notes and the master plan.

\begin{figure}[p]
\centering
\includegraphics[width=\textwidth]{fig_synthesis_4panel}
\caption{\small The ensemble on one page. \textbf{(a)}~The cost $J$ rises
monotonically with \kv{} across the members (dots); grey bands show the
internal-variability noise floor, and the star shows the reference sitting
above the members' trend --- the adjustment effect of Finding~1.
\textbf{(b)}~Spatial-structure survival of the accumulated temperature
sensitivity: $\approx$0.8 for the stable members, 0.49 for the degraded
$16\times$, $\approx$0 for the unstable members (open circles).
\textbf{(c)}~Amplitude of the sensitivity field as the adjoint works backward
(note the logarithmic axis spanning $\sim$40 decades): stable members hug the
reference (dashed black); unstable members leave it at the circled departure
points. \textbf{(d)}~Measured change in $J$ per member (solid) against the
adjoint's linear prediction (hatched); ``sign~\texttimes'' marks the four
members where even the sign is wrong.}
\label{fig:synth}
\end{figure}

\hd{Companion animations.}
Four self-contained interactive animations accompany this brief as separate
HTML files (each opens in any browser, no installation; time slider, play
control, and --- for the first two --- a depth selector):
\code{ADJtheta\_REF\_multidepth.html} and \code{ADJdiffkr\_REF\_multidepth.html}
(the reference sensitivities evolving at three depths),
\code{ADJtheta\_REF\_vs\_M3\_diff.html} (reference vs.\ the $2\times$ member,
side by side with their difference), and
\code{ADJtheta\_M5\_blowup\_log10.html} (the instability of
Figure~\ref{fig:blowup} frame by frame).
